\section{轻区与A80的形状混合}\label{sec:4-2}

在第~4.1~节所确立的跨区比较判据中，首要原则是将抽象的代数对称性与微观核子占据和具体的实验室可观测量严格锚定。沿核素图向中等质量区推进，$A\approx 60$ 与 $A\approx 80$ 区域（主要涵盖 $\mathrm{Ni}$、$\mathrm{Ge}$、$\mathrm{Se}$ 与 $\mathrm{Kr}$ 等同位素）展现出全核素图谱中极为剧烈、演化迅速的形状共存与构型混合现象。在这些核中，低激发谱上频繁出现反常的低激发 $0^+_2$ 态、强烈交错的能带结构以及受阻或增强的电磁跃迁。

然而，实验上观察到低激发 $0^+$ 态或非谐振动能级，绝不能简单断言体系必定存在多个具有深势垒隔离的宏观几何极小（即经典意义上的多形状共存）。本节将从两态混合的微观代数出发，依托 $^{68}\mathrm{Ni}$、$^{64}\mathrm{Ge}$、$^{74,76}\mathrm{Kr}$ 与 $^{76}\mathrm{Ge}$ 等经典实验原件，剖析电单极跃迁强度 $\rho^2(E0)$、静四极矩 $Q_s$ 的测量条件与 Kumar-Cline 四极几何不变量在检验模型中的实际角色，并明确与奇奇核手征-反射多重带的分工接口。

\subsection{两态与多态混合模型：能谱退化与物理可观测量}\label{sec:4-2-1}

理解形状共存与构型混合的最简洁物理起点是两态混合模型。设多体希尔伯特空间中存在两个未微扰的正交基矢 $|0^+_A\rangle$ 与 $|0^+_B\rangle$，它们分别对应具有不同内禀形变几何（例如一者为近球形或长椭球，另一者为扁椭球或大形变组态）的非微扰构型。设未微扰本征能量满足 $E_A \le E_B$，残余相互作用在两个构型之间诱发常数非对角耦合矩阵元 $V$（约定 $V > 0$）。

包含耦合的多体哈密顿量在基矢 $\{|0^+_A\rangle, |0^+_B\rangle\}$ 下写为：
\begin{equation}\label{eq:4-2-two-state-H}
\mathcal{H} = \begin{pmatrix} E_A & V \\ V & E_B \end{pmatrix}.
\end{equation}
对角化该哈密顿矩阵，物理观测到的两个正交定态 $|0^+_1\rangle$ 与 $|0^+_2\rangle$ 可表述为：
\begin{equation}\label{eq:4-2-two-state-states}
\begin{aligned}
|0^+_1\rangle &= \cos\theta_0 |0^+_A\rangle + \sin\theta_0 |0^+_B\rangle, \\
|0^+_2\rangle &= -\sin\theta_0 |0^+_A\rangle + \cos\theta_0 |0^+_B\rangle.
\end{aligned}
\end{equation}
要求物理定态之间非对角能量矩阵元严格消失，即
\begin{equation*}
\langle 0^+_1 | \mathcal{H} | 0^+_2 \rangle = \frac{1}{2}(E_B - E_A)\sin 2\theta_0 + V\cos 2\theta_0 = 0 ,
\end{equation*}
由此给出混合角与未微扰能隙的精确关系：
\begin{equation}\label{eq:4-2-two-state-mixing}
\Delta E = \sqrt{(E_B - E_A)^2 + 4V^2}, \qquad \tan 2\theta_0 = -\frac{2V}{E_B - E_A}.
\end{equation}
在 $E_A<E_B$、$V>0$ 时取 $-\pi/4<\theta_0<0$；当 $E_A=E_B$ 且 $V>0$ 时，连续极限为 $\theta_0=-\pi/4$。若能量相同且 $V=0$，二态空间简并，单凭哈密顿量不能固定混合角。

给定观测能隙 $\Delta E$，未扰动差 $(E_B-E_A)$ 与 $V$ 仍有多组解。混合大小由相对比例 $2V/(E_B-E_A)$ 决定：即使 $V$ 的绝对值很小，只要未扰动差更小，混合也可很强。二态能量拟合还不说明这两个基底是否对应彼此分离的形状极小；确定几何解释需要更多观测与模型条件。

电单极跃迁（Electric Monopole, $E0$）提供另一类约束。这里取点质子电荷算符 $\hat T(E0)=e\sum_{p=1}^{Z}r_p^2$，坐标相对核质心；每质子平均均方半径为 $\langle r^2\rangle_{A/B}=Z^{-1}\langle 0^+_{A/B}|\sum_p r_p^2|0^+_{A/B}\rangle$。除基底正交外，再假设该基底的非对角单极矩阵元可忽略；后一个假设不由正交性自动推出。由此得到：
\begin{equation}\label{eq:4-2-e0-matrix}
\langle 0^+_2 | \hat{T}(E0) | 0^+_1 \rangle = eZ\sin\theta_0\cos\theta_0\left(\langle r^2\rangle_B-\langle r^2\rangle_A\right).
\end{equation}
两个物理态之间的约化单极跃迁强度 $\rho^2(E0)$ 定义为无量纲比值：
\begin{equation}\label{eq:4-2-rho-e0-diff}
\rho^2(E0; 0^+_2 \to 0^+_1) \equiv \frac{|\langle 0^+_2 | \hat{T}(E0) | 0^+_1 \rangle|^2}{e^2 R_0^4} = \frac{1}{4} Z^2 \sin^2(2\theta_0) \frac{\left[ \Delta\langle r^2 \rangle_{\mathrm{intr}} \right]^2}{R_0^4},
\end{equation}
式中 $R_0=1.2A^{1/3}\,\mathrm{fm}$ 是所用参考半径。这个二态近似的 $\rho^2(E0)$ 约束混合因子与半径差平方的乘积：非零值需要二者均非零，给定强度却不能分别唯一确定混合角与半径差。半径差可与形变变化相关，具体关联仍须由构型模型检验。

\subsection{68Ni与64Ge：次闭壳破缺、寿命测定与形状证据边界}\label{sec:4-2-2}

在 $A\approx 60\text{--}70$ 区域，$^{68}\mathrm{Ni}$ 与 $^{64}\mathrm{Ge}$ 构成了理解微观壳层破缺与形状共存判定界限的典型案例。

$^{68}\mathrm{Ni}$（$Z=28, N=40$）质子处于闭壳，中子填满 $1f_{5/2}, 2p_{3/2}, 2p_{1/2}$ 轨道，其费米面恰位于 $N=40$ 次闭壳。早期谱学曾根据较高的第一激发 $2^+_1$ 态能量（2033 keV）将其视作良好球形核。然而后续实验逐渐揭示出复杂的低能激发现象。Crider 等（2016）\cite{crider-2016-68ni-shape-coexistence}利用 $\beta$ 衰变谱学结合超快计时技术，对 $^{68}\mathrm{Ni}$ 的高能态与低激发态进行了寿命与分支比测量。核对 Crider 等（2016）原件 Table 1：
\begin{itemize}
\item 第一激发 $0^+_2$ 态位于 1605 keV，属于长寿命异能态（文献采用半衰期 $T_{1/2} = 268(12)\,\mathrm{ns}$），其向基态通过纯 $E0$ 内转换电子退激，采用强度为 $\rho^2(E0; 0^+_2 \to 0^+_1) \times 10^3 = 7.6(4)$；
\item Crider 等（2016）在该工作中重点测定了第二激发 $0^+_3$ 态（2511 keV）的寿命，测得其半衰期为亚纳秒量级的 $T_{1/2} = 0.57(5)\,\mathrm{ns}$；
\item 该 $0_3^+$ 态到 $2_1^+$ 的分支大于0.981，相应 $B(E2)=39.0(34)\,e^2\mathrm{fm}^4$。Table 1对其到基态及 $0_2^+$ 的E0强度仅给上限，分别为 $\rho^2(E0)<0.0050$ 和 $<0.0258$；这些上限不是测得的非零强度。
\item 同文对 $^{70}\mathrm{Ni}$ 候选 $(0_2^+)$ 采用或测定的量须单独列出：半衰期为 $1.65^{+0.30}_{-0.25}\,\mathrm{ns}$，到 $2_1^+$ 的E2分支大于0.66，而到基态的E0仅给上限。不能将其与 $^{68}\mathrm{Ni}$ 的另一态合成一次测量。
\end{itemize}
作者比较LNPS壳模型与采用A3DA相互作用的MCSM，将 $^{68}\mathrm{Ni}$ 的 $0_1^+$、$0_2^+$、$0_3^+$ 分别与球形、扁椭和长椭构型联系起来。这是模型对态结构的解释。原文第4页还列出 $0_3^+$ 的半衰期预言：A3DA与LNPS分别约108 ns和1.5 ns，与实测半衰期0.57(5) ns不同。能级和形状标签相近的模型仍可给出差异很大的绝对率，寿命因而提供另一项检验，而不能单独当作形状的直接成像。

在自共轭核 $^{64}\mathrm{Ge}$（$N=Z=32$）中，证据边界同样需要严格澄清。Starosta 等（2007）\cite{doctoral-r1-g26-starosta-2007-64ge-rdm-e2}利用单中子敲出反应 $^{9}\mathrm{Be}(^{65}\mathrm{Ge}, ^{64}\mathrm{Ge}+\gamma)X$ 结合反冲距离多普勒频移法（RDM），系统测定了低自旋态的寿命。核对原件 PDF 第 3 页 Table I 及表下说明：
\begin{itemize}
\item 第一激发 $2^+_1$ 态（902 keV）的实测平均寿命为 $\tau(2^+_1) = 3.3(5)\,\mathrm{ps}$，向下跃迁几率为 $B(E2; 2^+_1 \to 0^+_1) = 410(60)\,e^2\mathrm{fm}^4$；
\item 第二激发 $2^+_2$ 态（1579 keV）的实测平均寿命为 $\tau(2^+_2) = 8^{+4}_{-2}\,\mathrm{ps}$，向下跃迁几率分别为 $B(E2; 2^+_2 \to 2^+_1) = 620(210)\,e^2\mathrm{fm}^4$ 以及强烈受阻的交叉跃迁 $B(E2; 2^+_2 \to 0^+_1) = 1.5(5)\,e^2\mathrm{fm}^4$；
\item 原文Table I把静四极矩 $Q(2_1^+)$、$Q(2_2^+)$ 列在GXPF1A理论计算栏；这些数值不能写成该次寿命实验直接测得的静矩。
\item 原文作者据此给出的物理结论是``$^{64}\mathrm{Ge}$ 呈现集体 $\gamma$ 软非谐振动特征''。
\end{itemize}
因此，强 $B(E2)$ 证明了核内存在显著的四极集体性，但在缺乏对角静矩实测与高阶不变量的情况下，绝不能把理论预测列当成实验实测值，更不能将该核直接当作已经实验证实的多形状共存。

\subsection{74,76,78Kr：多步库仑激发、能带混合与静四极矩测量}\label{sec:4-2-3}

在氪（$\mathrm{Kr}$）同位素链中，随着中子数向中子亏损一侧推进，势能面表现出强烈的长椭与扁椭竞争。由于实验室系中 $0^+$ 基态由角动量守恒定理决定其一阶静四极矩严格为零（$\langle 0^+ | \hat{Q}_{20} | 0^+ \rangle \equiv 0$），内禀电荷各向异性的对角信息在低自旋下主要体现于带内激发态（如 $2^+_1, 2^+_2$）的静四极矩 $Q_s$ 中。

根据第2.1节限定于轴对称 $K=0$ 带的几何关系，对于纯轴对称长椭构型（$Q_0 > 0$），基态带 $2^+_1$ 态的静四极矩预期为负值（$Q_s(2^+_1) < 0$）；而对于扁椭构型（$Q_0 < 0$），则预期为正值（$Q_s(2^+_1) > 0$）。然而在实际原子核中，由于构型混合的存在，观测到的定态静矩是各基矢分量矩阵元的相干叠加。

Clément 等（2007）\cite{doctoral-r1-g27-clement-2007-prc-krypton-coulex}在 GANIL/SPIRAL 装置上，利用放射性束 $^{74}\mathrm{Kr}$ 与稳定束 $^{76}\mathrm{Kr}$ 轰击 $^{208}\mathrm{Pb}$ 靶进行了多步安全库仑激发，通过 GOSIA 程序拟合求解了包含数十个能级矩阵元的网络。查阅 Clément 等（2007）原件表格：
\begin{itemize}
\item 在 $^{76}\mathrm{Kr}$ 中，原件 PDF 第 11 页 Table VI 列出的实验结果为：$2^+_1$ 态的静四极矩为 $Q_s(2^+_1) = -0.7(2)\,e\mathrm{b}$（约化矩阵元 $\langle 2^+_1 || M(E2) || 2^+_1 \rangle = -0.9(3)\,e\mathrm{b}$）；
\item 在更缺中子的 $^{74}\mathrm{Kr}$ 中，原件 PDF 第 15 页 Table XI 列出的实验拟合结果为：$2^+_1$ 态的静四极矩为 $Q_s(2^+_1) = -0.53^{+0.24}_{-0.23}\,e\mathrm{b}$（约化矩阵元 $\langle 2^+_1 || M(E2) || 2^+_1 \rangle = -0.70^{+0.31}_{-0.30}\,e\mathrm{b}$）。
\end{itemize}
由此可见，实验直接测得的 $^{74}\mathrm{Kr}$ 与 $^{76}\mathrm{Kr}$ 第一激发 $2^+_1$ 态的静四极矩测量中心值\textbf{均为负值}，并非网络文献或二手笔记中所误传的一正一负数值。

那么 Clément 等是如何确立两核中存在长椭与扁椭形状共存的？其论证并非建立在单个静矩的孤立符号上，而是基于两能带混合模型（two-band mixing model）对全能带跃迁矩阵元与能级的全局自洽拟合：
\begin{itemize}
\item 在未微扰表象中，假定存在一个长椭带和一个扁椭带；
\item 全局拟合表明，在 $^{76}\mathrm{Kr}$ 中，未微扰基态带由长椭构型主导，长椭构型处于能量低端；
\item 而在 $^{74}\mathrm{Kr}$ 中，未微扰能隙发生反转，未微扰扁椭构型的带头能量低于长椭构型带头；通过强烈的带间混合矩阵元（$V \approx 330\,\mathrm{keV}$），混合后的物理基态具有显著的未微扰扁椭成分（约占 $50\%\text{--}60\%$），而物理激发带则与长椭构型强烈交织。
\end{itemize}
这一严密的原始证据表明：形状反转与共存的判定依赖于包含多个态跃迁与静矩的完整混合模型反演，绝不能用对角单态静矩的简单符号去代替能带波函数的复杂混合。

\subsection{72,74,76Ge：刚性三轴 vs 软势竞争与Kumar-Cline不变量提取}\label{sec:4-2-4}

在锗（$\mathrm{Ge}$）同位素链中，物理探讨的核心聚焦于势能面的三轴刚性与 $\gamma$ 软性。

Ayangeakaa 等（2019）\cite{ayangeakaa-2019-rigid-triaxial-ge76}利用 GRETINA 阵列结合 CHICO2 探测器，对 $^{76}\mathrm{Ge}$ 实施了多步库仑激发，在有限能级范围内拟合了电磁矩阵元网络，并结合采用的寿命、分支等约束。该工作采用 Kumar-Cline 几何不变量方法（见式~\eqref{eq:4-1-q2-invariant}与式~\eqref{eq:4-1-q3-invariant}），提取了四极变形不变量 $\langle Q^2 \rangle$ 以及三轴不变量 $\langle Q^3 \cos 3\delta \rangle$。核对原件 PDF 第 5–6 页：该工作在 Fig.~4 与 Fig.~5 中分析了基态带与 $\gamma$ 带的等效三轴平均角（约处于 $25^\circ\text{--}27^\circ$ 区间），并评估了高阶不变量分布的统计弥散（如 $\sigma_{\langle Q^2 \rangle}$ 与 $\sigma_{\langle \cos 3\delta \rangle}$）。原作者据此结合模型指出，$^{76}\mathrm{Ge}$ 的低自旋态展现出较为稳定的刚性三轴形变特征。

然而，正如第~2.4~节所述，将几何不变量与刚性转子模型对接时必须牢记其适用边界：不变量求和规则在实际应用中受制于有限的能级截断与不可测的高能矩阵元；分散度反映的是电荷分布高阶矩的统计涨落，不能将其设定为某种机械门槛去绝对宣称消除了全部软性。不变量的核心价值在于提供与微观多体波函数直接比对的模型无关坐标，证明 $^{76}\mathrm{Ge}$ 的势能面在三轴非零处具有明确的局域结构。

与 $^{76}\mathrm{Ge}$ 相比，同位素链较轻端的 $^{72}\mathrm{Ge}$ 展现出截然不同的低能共存行为。在 $^{72}\mathrm{Ge}$ 中，第一激发态是位于 691 keV 的 $0^+_2$ 态，比第一 $2^+_1$ 态（834 keV）还要低。该态主要通过发射内转换电子的强 $E0$ 跃迁向基态衰变，这是源于质子中子跨次壳激发诱导的构型共存典型表象。从 $^{72}\mathrm{Ge}$ 的低能形状共存，演化至 $^{76}\mathrm{Ge}$ 展现出的稳定三轴特征，体现了中子填充如何平滑局域壳隙并重塑核芯四极动力学。

\subsection{A≈60奇奇核手征探索与残余相互作用：Ding 2025微观模型分析}\label{sec:4-2-5}

当分析从偶偶核延伸至奇奇核时，未配对准粒子与非轴对称核芯的耦合引出了核手征性的探索。Ding 等（2025）\cite{ding-2025-chiral-rotation-a60-residual-interactions}在理论上开展了一项微观粒子转子模型（PRM）研究，探索了 $A\approx 60$ 质量区奇奇核中手征对称性破缺的可能性。

需要明确该论文的理论定位：该工作是一项基于协变密度泛函（CDFT）微观输入构建的理想化 PRM 理论推演，采用 $\pi f_{7/2}^{-1} \otimes \nu g_{9/2}$ 准粒子构型耦合至具有 $\beta=0.25, \gamma=34^\circ$ 的三轴核芯，\textbf{并非报告新的实验测定或在钴同位素中发现了物理实测能带}。该工作核心关注的是轻区核中残余质中相互作用 $V_{pn}$ 的动力学角色。

在奇奇核转动带中，Ding 等考察了能量交错参数 $S(I)$ 的演化：
\begin{equation}\label{eq:4-2-si-staggering}
S(I) = \frac{E(I) - E(I-1)}{2I}.
\end{equation}
在此必须在物理概念上做出严格区分：
\begin{enumerate}
\item 式~\eqref{eq:4-2-si-staggering}所定义的 $S(I)$，是奇奇核单个转动带内相邻自旋能级间隔除以总角动量的量，反映的是有效转动能谱随自旋增加的单调性或奇偶微细起伏；
\item 偶偶核 $\gamma$ 振动带中常引用的交错参数则定义为二阶差分比值：$S(I) = \frac{[E(I) - 2E(I-1) + E(I-2)]}{E(2^+_1)}$，用以探测能级是否按对成簇聚集（如 Wilets-Jean 软核中 $3^+-4^+$、$5^+-6^+$ 的简并结构）。
\end{enumerate}
两处虽然共用字母 $S(I)$，但其数学阶数与物理目标截然相异。绝不能将奇奇核中的能量间隔 $S(I)$ 随残余相互作用的变动，直接当作偶偶核刚性或软性 $\gamma$ 判据的反例。

Ding 等（2025）的实际理论贡献在于指出：在价核子数目较少的轻区体系中，残余质中相互作用 $V_{pn}$ 会显著重整化准粒子角动量在三个主轴上的投影几率与取向角动量几何，使得带内能量间隔 $S(I)$ 的奇偶振荡幅度甚至符号发生显著改变。这一理论结果提醒我们：在利用唯象几何模型拟合奇奇核谱形时，必须谨慎评估微观残余多极力的重整化效应，避免将剩余力引起的能级移动简单归咎于几何形变参数的改变。

\subsection{与奇奇核（74As与78Br）的物理接口界定}\label{sec:4-2-6}

最后，本节明确界定与后续第~4.5~节的科学分工接口。在 $A\approx 80$ 奇奇核中，$^{74}\mathrm{As}$ 与 $^{78}\mathrm{Br}$ 是探索手征几何与可能存在的内禀反射不对称性的前沿对象。

从低能轨道占据来看，$^{78}\mathrm{Br}$（$Z=35, N=43$）与 $^{74}\mathrm{As}$（$Z=33, N=41$）的奇价质子与奇价中子分别填充在负宇称 $fp$ 壳层与正宇称 $1g_{9/2}$ 轨道上。价核子的极化效应对原本处于形状过渡区的偶偶核芯施加了三轴与八极多极驱动力。然而，正如刘晨在其博士学位论文\cite{liu-chen-2016-78br-chirality-reflection-symmetry-thesis}中针对 $^{78}\mathrm{Br}$ 高自旋在束实验分析时所客观陈述的（见其论文第 90 页总结与展望）：
\begin{quote}
``电磁跃迁几率 $B(M1)/B(E2)$ 比值也存在着奇偶摇摆效应。这些特征都与手征带的实验特征相吻合，但由于目前的实验数据较为有限，无法得到准确的结论。这些核中的双带是否为手征带，仍然需要实验上进一步的研究。''
\end{quote}

本节严格恪守这一客观证据边界：本节仅在低能形状混合与轨道占据的层次上，将 $^{74}\mathrm{As}$ 与 $^{78}\mathrm{Br}$ 作为奇核子极化核芯形状的接口输入；关于在更高自旋下是否真实形成了四条近简并的手征-反射联立破缺多重带，以及微观多极 $E1$ 连接跃迁与 TAC-CDFT 模型的详细检验，将完全交由第~4.5~节进行全方位的独立审视与深入论证，本节不作越界的确定性推断。
